% 使用 XeLaTeX 编译；中文采用 TeX Live 自带的 Fandol 字体。
\documentclass[11pt,a4paper,fontset=fandol]{ctexart}
\usepackage[margin=2.1cm]{geometry}
\usepackage{amsmath,amssymb,mathtools}
\usepackage{fancyhdr,needspace}
\usepackage[hidelinks]{hyperref}
\hypersetup{pdftitle={线性代数第五次作业},pdfsubject={MATH1205：逆矩阵、伴随矩阵与矩阵幂}}
\setlength{\parindent}{0pt}
\setlength{\parskip}{0.25em}
\setlength{\headheight}{15pt}
\renewcommand{\arraystretch}{1.08}
\pagestyle{fancy}
\fancyhf{}
\fancyhead[L]{MATH1205\quad 线性代数}
\fancyhead[R]{第五次作业}
\fancyfoot[C]{\thepage}
\renewcommand{\headrulewidth}{0.3pt}
\ctexset{section={format=\large\bfseries,beforeskip=1.2ex,afterskip=0.7ex},
subsection={format=\normalsize\bfseries,beforeskip=0.9ex,afterskip=0.5ex}}
\newcommand{\ans}[1]{\boxed{#1}}
\newcommand{\E}{E}
\newcommand{\problem}[1]{\Needspace{9\baselineskip}\section*{#1}}

\begin{document}
\begin{center}
{\LARGE\bfseries 线性代数第五次作业}\par
\vspace{0.3em}
{\small MATH1205}
\end{center}
\vspace{-0.5em}
以下 $E$ 表示相应阶数的单位矩阵，$O$ 表示零矩阵，$A^*$ 表示 $A$ 的伴随矩阵。
题号沿用原题。

\problem{第 24 题：解含伴随矩阵的矩阵方程}
已知
\[
A^*BA=2BA-8E,\qquad
A=\begin{pmatrix}1&2&-2\\0&-2&4\\0&0&1\end{pmatrix}.
\]
因为 $|A|=1\cdot(-2)\cdot1=-2\ne0$，所以 $A$ 可逆，且 $AA^*=-2E$。
在原式两端左乘 $A$，得
\[
-2BA=2ABA-8A.
\]
再在两端右乘 $A^{-1}$，整理得
\[
-2B=2AB-8E,\qquad (A+E)B=4E.
\]
而
\[
A+E=\begin{pmatrix}2&2&-2\\0&-1&4\\0&0&2\end{pmatrix},
\qquad |A+E|=-4\ne0.
\]
对 $(A+E)B=4E$ 的各列从最后一行向上回代，得到原方程的唯一解
\[
\ans{B=4(A+E)^{-1}=
\begin{pmatrix}2&4&-6\\0&-4&8\\0&0&2\end{pmatrix}}.
\]
上述各步均为等价变形，且 $(A+E)B=4E$，故所得 $B$ 满足原方程。

\problem{第 25 题：求逆矩阵}
\subsection*{(2)}
记
\[
R=\begin{pmatrix}\cos\theta&-\sin\theta\\\sin\theta&\cos\theta\end{pmatrix}.
\]
由 $\cos^2\theta+\sin^2\theta=1$，直接相乘可得
\[
R\begin{pmatrix}\cos\theta&\sin\theta\\-\sin\theta&\cos\theta\end{pmatrix}
=\begin{pmatrix}1&0\\0&1\end{pmatrix}.
\]
故对任意实数 $\theta$，都有
\[
\ans{R^{-1}=\begin{pmatrix}\cos\theta&\sin\theta\\-\sin\theta&\cos\theta\end{pmatrix}}.
\]

\newpage
\section*{第 25 题（续）}
\subsection*{(4)}
设
\[
U=\begin{pmatrix}
1&a&a^2&a^3\\0&1&a&a^2\\0&0&1&a\\0&0&0&1
\end{pmatrix},\qquad
V=\begin{pmatrix}
1&-a&0&0\\0&1&-a&0\\0&0&1&-a\\0&0&0&1
\end{pmatrix}.
\]
记 $U$ 的各列为 $C_1,C_2,C_3,C_4$，标准单位列向量为 $e_1,\ldots,e_4$。
则
\[
UV=(C_1,\ C_2-aC_1,\ C_3-aC_2,\ C_4-aC_3)
=(e_1,e_2,e_3,e_4)=E.
\]
所以 $\ans{U^{-1}=V}$。此式在题设 $a\ne0$ 时成立；实际上 $a=0$ 时同样成立。

\problem{第 28 题：利用相似关系计算矩阵幂}
已知
\[
A=\begin{pmatrix}1&0&0\\2&-1&0\\2&1&1\end{pmatrix},\qquad
B=\begin{pmatrix}1&0&0\\0&0&0\\0&0&-1\end{pmatrix},
\qquad XA=AB.
\]
\subsection*{(1) 证明 $X^n=AB^nA^{-1}$}
由于 $|A|=-1\ne0$，在 $XA=AB$ 两端右乘 $A^{-1}$，得 $X=ABA^{-1}$。
对正整数 $n$ 用数学归纳法。$n=1$ 时结论成立。
若 $X^m=AB^mA^{-1}$，则
\[
X^{m+1}=X^mX
=AB^mA^{-1}ABA^{-1}
=AB^m(A^{-1}A)BA^{-1}
=AB^{m+1}A^{-1}.
\]
故
\[
\ans{X^n=AB^nA^{-1}\qquad(n\ge1)}.
\]

\Needspace{14\baselineskip}
\subsection*{(2) 求 $X^5$}
因为 $B$ 为对角矩阵，
\[
B^5=\operatorname{diag}(1^5,0^5,(-1)^5)
=\operatorname{diag}(1,0,-1)=B.
\]
又由直接计算可得
\[
A^{-1}=\begin{pmatrix}1&0&0\\2&-1&0\\-4&1&1\end{pmatrix},
\qquad
AB=\begin{pmatrix}1&0&0\\2&0&0\\2&0&-1\end{pmatrix}.
\]
因此
\[
\ans{X^5=ABA^{-1}
=\begin{pmatrix}1&0&0\\2&0&0\\2&0&-1\end{pmatrix}
 \begin{pmatrix}1&0&0\\2&-1&0\\-4&1&1\end{pmatrix}
=\begin{pmatrix}1&0&0\\2&0&0\\6&-1&-1\end{pmatrix}}.
\]

\newpage
\problem{第 29 题：由逆矩阵求原矩阵及其伴随矩阵}
记
\[
C=A^{-1}=\begin{pmatrix}3&-1&1\\1&1&0\\2&1&1\end{pmatrix}.
\]
按第一行展开得
\[
|C|=3(1\cdot1-0\cdot1)
-(-1)(1\cdot1-0\cdot2)+(1\cdot1-1\cdot2)=3.
\]
计算 $C$ 的各代数余子式并转置，利用 $C^{-1}=C^*/|C|$，得
\[
\ans{A=C^{-1}=\frac13
\begin{pmatrix}1&2&-1\\-1&1&1\\-1&-5&4\end{pmatrix}},
\qquad |A|=\frac1{|C|}=\frac13.
\]
利用 $A^*=|A|A^{-1}$，得
\[
\ans{A^*=\frac13
\begin{pmatrix}3&-1&1\\1&1&0\\2&1&1\end{pmatrix}}.
\]
由于 $A^*=\frac13 C$，所以
\[
\ans{(A^*)^{-1}=3C^{-1}=3A
=\begin{pmatrix}1&2&-1\\-1&1&1\\-1&-5&4\end{pmatrix}}.
\]
又因 $A^*$ 为三阶矩阵，
\[
|A^*|=\left|\frac13 C\right|
=\left(\frac13\right)^3|C|=\frac19.
\]
从而
\[
\ans{(A^*)^*=|A^*|(A^*)^{-1}
=\frac19\begin{pmatrix}1&2&-1\\-1&1&1\\-1&-5&4\end{pmatrix}}.
\]

\problem{第 31 题：证明 $A-E$ 可逆并求其逆}
已知 $A^2=2A$，于是
\[
(A-E)(A-E)=A^2-2A+E=2A-2A+E=E.
\]
根据逆矩阵的定义，$A-E$ 可逆，且
\[
\ans{(A-E)^{-1}=A-E}.
\]

\newpage
\problem{第 33 题：证明 $A^*$ 不可逆}
设 $A$ 为 $n$ 阶方阵，$n\ge2$，且 $A^2=2A$、$A\ne2E$。

先证明 $A$ 不可逆。若 $A$ 可逆，则在 $A^2=2A$ 两端左乘 $A^{-1}$，得到
$A=2E$，与条件矛盾。因此 $|A|=0$，由伴随矩阵恒等式可得
\[
AA^*=|A|E=O.
\]
若 $A^*$ 可逆，则在上式两端右乘 $(A^*)^{-1}$，得到 $A=O$。
此时因 $n\ge2$，$A$ 的每个余子式都是至少一阶的零矩阵的行列式，均为零，
从而 $A^*=O$，与 $A^*$ 可逆矛盾。因此
\[
\ans{A^*\text{ 不可逆}}.
\]

\problem{第 35 题：幂零矩阵的有限几何级数}
已知 $A^k=O$，其中 $k$ 为正整数。令
\[
S=E+A+A^2+\cdots+A^{k-1}.
\]
直接计算得
\begin{align*}
(E-A)S
&=(E+A+\cdots+A^{k-1})-(A+A^2+\cdots+A^k)\\
&=E-A^k=E,\\
S(E-A)
&=(E+A+\cdots+A^{k-1})-(A+A^2+\cdots+A^k)\\
&=E-A^k=E.
\end{align*}
所以 $E-A$ 可逆，且
\[
\ans{(E-A)^{-1}=E+A+A^2+\cdots+A^{k-1}}.
\]
当 $k=1$ 时，$A=O$，右端的和仅有 $E$ 一项，结论仍成立。

\problem{第 38 题：计算伴随矩阵及矩阵组合的行列式}
已知 $A$ 是五阶矩阵，且 $|A|=\frac12\ne0$，故 $A$ 可逆，
\[
A^*=|A|A^{-1}=\frac12A^{-1},\qquad |A^{-1}|=2.
\]
因此
\[
\ans{|A^*|=\left|\frac12A^{-1}\right|
=\left(\frac12\right)^5\cdot2=\frac1{16}}.
\]
另一方面，
\[
(2A)^{-1}+3A^*
=\frac12A^{-1}+3\cdot\frac12A^{-1}
=2A^{-1}.
\]
再取行列式，得
\[
\ans{\left|(2A)^{-1}+3A^*\right|
=|2A^{-1}|=2^5|A^{-1}|=64}.
\]
\end{document}
